\documentclass[18pt]{article}
\usepackage{geometry}
\geometry{top=3cm}
\usepackage{amsmath,amsfonts,amssymb}
\usepackage{natbib}
\usepackage{booktabs}
\usepackage[utf8]{inputenc}
\usepackage[T1]{fontenc}
\usepackage{amssymb}
\usepackage[version=4]{mhchem}
\usepackage{stmaryrd}
\usepackage{bbold}
\usepackage{nopageno}
\usepackage{hyperref}
\usepackage{bookmark}
\usepackage{relsize}
\author{Adam Handwerger}
\pagestyle{plain}

\begin{document}
\begin{flushleft}
2 Optimal Ordering\\
\vspace{.4cm}
Let $x_i \in \mathcal{X}$ and $y_i \in \{-1,+1\} \hspace{.3cm} 1 \leq i\leq n $. Define\\
\vspace{.4cm}
\hspace{3cm} $I_-=\{i | y_i=-1\}$ with $n_{-}=$ num. of elements in $I_-$\\
\hspace{3cm} $I_+=\{i | y_i=+1\}$ with $n_{+}=$ num. of elements in $I_+$\\
\vspace{.4cm}
And let $n=n_-+n_+$. We seek an f which minimizes the functional\\
\vspace{.4cm}
\hspace{2cm} $J(f)=\dfrac{1}{n_- n_+} \sum\limits_{i \in I_-}\sum\limits_{j \in I_+}(1-f(f(x_j)-f(x_i)))+\lambda \Vert f \Vert^2_H$\\
\vspace{.4cm}
1. Justify that such an $f^*$ can be expressed a function indexed by $\alpha$ of the form\\
\vspace{.4cm}
\hspace{4cm} $f_\alpha^*(x)=\sum\limits_{i=1}^n \alpha_i K(x,x_i)$\\
\vspace{.4cm}
Let H as a Hilbert space and let $f\in H.$ Since H is a Hilbert space it can be expressed as\\
\vspace{.4cm}
\hspace{4cm} $H=H_\mathcal{V}\oplus H_{\perp}$\\
\vspace{.4cm}
where both $H_\mathcal{V}$ and $H_{\perp}$ are orthogonal subspaces. This means any $f \in H$ can expressed as\\
\vspace{.4cm}
\hspace{4cm} $f^*=f_\mathcal{V}+f_{\perp}$.\\
\vspace{.4cm}
We contruct another subspace of H defined as\\
\vspace{.4cm}
\hspace{3cm} $H_\mathcal{V} =\left\{ f : f(x)=\sum \alpha_i K(x,x_i) \right\}$\\
\vspace{.4cm}
By construction $K(\cdot,x_i) \in  H_\mathcal{V}$, therefore for any $1 \leq i \leq n$  we have\\
\vspace{.4cm}
\hspace{1cm} $f(x_i)=\langle f_\mathcal{V}+f_\perp,K(\cdot,x_i) \rangle=\langle f_\mathcal{V}, K(\cdot,x_i)\rangle+\langle f_\perp, K(\cdot,x_i)\rangle=f_\mathcal{V}(x_i)$.\\
\vspace{.4cm}
Therefore since $f^* \in H$ is the f that minimizes $J(f)$, then we have\\
\vspace{.4cm}
\hspace{2cm} $J(f^*)-J(f_\mathcal{V})=\lambda(\vert| f_\mathcal{V}\vert|^2+\vert|f_{\perp}\vert|^2)-\lambda\vert| f_\mathcal{V}\vert|^2=\lambda\vert| f_\mathcal{V}\vert|^2 \geq 0$\\
\vspace{.4cm}
Other the other hand $f^*$ minimizes $J(f)$ and so\\
\vspace{.4cm}
\hspace{4cm} $J(f^*)-J(f_\mathcal{V}) \leq 0$\\
\vspace{.4cm}
Both equations together implies $\vert|f_{\perp}\vert|=0$ which futher implies that $f_{\perp}=0$.\\
\vspace{.4cm}
Hence any $f^*\notin H_{\perp}$ and so $f^*=f_\mathcal{V}$. I.e all f which minimize $J(f)$ must be in $H_\mathcal{V}$\\
\vspace{.4cm}
2. Write J(f) in terms of $\alpha$.\\
\vspace{.4cm}
Substituting the form of $f \in \mathcal{V}$ into $J(f)$ gives\\
\vspace{.4cm}
$J(f)=\dfrac{1}{n_-n_+} \sum\limits_{ij} \left(1-\sum\limits_k \alpha_k K(x_k,x_j)-\sum\limits_k \alpha_k K(x_k,x_i)\right)+\lambda\alpha^{\top}K \alpha$=\\
$\dfrac{1}{n_-n_+} \sum\limits_{ij} \left(1-\sum\limits_k \alpha_k([K]_{kj}-[K]_{ki})\right)+\lambda\alpha^{\top}K \alpha=\frac{1}{n_-n_+} \sum\limits_{ij}\left(1-\alpha^{\top}(K_j-K_i)\right)+\lambda\alpha^{\top}K \alpha=$\\
$1- \dfrac{\alpha^{\top}}{ n_{-} n_{+}} \sum\limits_{j \in I_+}\sum\limits_{i \in I_-}\left(K_{j}-K_{i} \right)+\lambda\alpha^{\top}K \alpha$\\
\vspace{.4cm}
with $K_j$ is defined as jth collumn of K and likewise for $K_i$.\\
3. Furthermore if we define $K_-=\dfrac{1}{n_-}\sum\limits_{i \in I_-} K_i$ and likewise, $K_+-=\dfrac{1}{n_+}\sum\limits_{j \in I_+} K_j$.\\
Substituting,\\
\vspace{.4cm}
$J(f)=1- \frac{\alpha^{\top}}{ n_{-} n_{+}} \sum\limits_{j \in I_+}\sum\limits_{i \in I_-}\left(K_{j}-K_{i} \right)+\lambda\alpha^{\top}K \alpha=1-\alpha^{\top}\left(\dfrac{K_+}{n_{-}}-\dfrac{K_-}{n_{+}}\right)+\lambda\alpha^{\top}K \alpha$\\
\vspace{.4cm}
*From discussion in class, that the $n\pm$ in the denominators vanish after summing over $i,j$.\\
\vspace{.4cm}
4. Calculate $\triangledown_\alpha J(\alpha)$.\\
\vspace{.4cm}
Using $\triangledown_\alpha \alpha^{\top}b=b$ and $\triangledown_\alpha^{\top} K \alpha=(K^\top+K) \alpha=2K\alpha$ gives\\
\vspace{.4cm}
\hspace{3cm} $\triangledown_\alpha J(f)=\left(\dfrac{K_+}{n_{-}}-\dfrac{K_-}{n_{+}}\right)+2K\lambda \alpha$\\
\vspace{.4cm}
5. Setting the gradient to zero immediately gives $\alpha^*$:\\
\vspace{.4cm}
\hspace{3cm} $\alpha^*=\dfrac{K^{-1}}{2 \lambda}\left(\dfrac{K_+}{n_{-}}-\dfrac{K_-}{n_{+}}\right)$\\
\vspace{.4cm}
6. Using the kernel, $K(x,x^{\prime})=x^{\top}x^{\prime}$ with associated RKHS\\
\vspace{.4cm}
\hspace{1cm} $H={f_v; f_v(x)=v^{\top}}x,$ with $v \in \mathbb{R}^d$ and $\langle f_v,f_w \rangle=v^{\top}w$\\
\vspace{.4cm}
verify the reproducing property of K.\\
\vspace{.4cm}
By inspection, $K(x,\cdot)=f_x \implies $\\
$\langle f_v,K(x,\cdot) \rangle=\langle f_v,f_x \rangle=v^{\top}x=f_v(x)$\\
\vspace{.4cm}
7. Compute $J(f(\alpha^*))$.\\
\vspace{.4cm}
Let $f(\alpha^*) \equiv f^*(x)=\sum\limits_k^n \alpha_k K(x,x_k)=\sum\limits_k^n \alpha_k^* x^{\top}x_k=$\\
$\sum\limits_k^n \alpha_k^* x_k^{\top}x=\beta^{\top}x$ with $\beta^{\top} \equiv \sum\limits_k \alpha_k^* x_k^{\top}$\\
\vspace{.4cm}
Hence\\
\vspace{.4cm}
$\beta^{\top}=\mathlarger\sum\limits_{k} \left[ \dfrac{K^{-1}}{2 \lambda}\left(\dfrac{K_+}{n_{-}}-\dfrac{K_-}{n_{+}}\right)\right]_k x_k^{\top}$=\\
$\mathlarger\sum\limits_{k} \left[ \dfrac{K^{-1}}{2 \lambda n_+ n_-}\left(\sum\limits_{j \in I_-} K_j-\sum\limits_{i \in I_+} K_i\right)\right]_k x^{\top}_k$\\
\vspace{.4cm}
Defining $x^+ \equiv \frac{1}{n_+} \sum_j^{n_+} x_j$ and $x^- \equiv \frac{1}{n_-} \sum_i^{n_-} x_i$ gives\\
$\sum\limits_{j \in I_+} K_j=$\\
\vspace{.4cm}
$$
    \begin{bmatrix}
        x_{1}^{\top}x_{j_1}\\x_{2}^{\top} x_{j_1}\\ \vdots \\ x_{n}^{\top} x_{j_1}
    \end{bmatrix}
    +\cdots+
    \begin{bmatrix}
        x_{1}^{\top} x_{j_{n_+}}\\x_{2}^{\top} x_{j_{n_+}}\\ \vdots \\ x_{n}^{\top} x_{j_{n_+}}
    \end{bmatrix}
    =
    \begin{bmatrix}
        x_{1}^{\top} \sum \limits_{j \in I_+}^{n_+} x_{j} \\ \vdots\\ x_{n}^{\top}\sum\limits_{j \in I_+}^{n_+} x_{j}
    \end{bmatrix} =
    \begin{bmatrix}
        x_1^{(1)} & x_1^{(2)}& \cdots &x_1^{(d)}\\
        \vdots    & \vdots    &        &\vdots\\
        x_n^{(1)} & x_n^{(2)}& \cdots &x_n^{(d)}\\
    \end{bmatrix}  \left(\dfrac{x^+}{n_-}\right)
$$
\vspace{.4cm}
and likewise for $\sum\limits_{i \in I+} K_i$ and therefore $\beta^{\top}$ can written as\\
\vspace{.4cm}
$\hspace{3cm} \beta^{\top}=\mathlarger\sum\limits_{k} \left[ \dfrac{K^{-1}}{2 \lambda}\left(\dfrac{K_+}{n_{-}}-\dfrac{K_-}{n_{+}}\right)\right]_k x_k^{\top}$=\\

\vspace{.4cm}
$$
\mathlarger\sum\limits_{k}^n \left[ \dfrac{K^{-1}}{2 \lambda}
\begin{bmatrix}
    x_1^{(1)} & x_1^{(2)}& \cdots &x_1^{(d)}\\
    \vdots    & \vdots    &        &\vdots\\
    x_n^{(1)} & x_n^{(2)}& \cdots &x_n^{(d)}
\end{bmatrix}  \left(\dfrac{x_+}{n_-}-\dfrac{x_-}{n_+}\right)\right]_k x_k^{\top}
$$\\
\vspace{.4cm}
To summerize the method to find $\beta^{\top}$:\\
\vspace{.4cm}
First, caluclate $x_+$, and $x_-$ from the data points to find $\left(\dfrac{x_+}{n_-}-\dfrac{x_-}{n_+}\right)$. This is $(d\times1)$ vector.\\
\vspace{.4cm}
Second, create matrix of data points with componenets in each row. This is a $(n \times d)$ matrix.\\
\vspace{.4cm}
Third, calulate the inverse of matrix K. This is an $(n \times n)$ matrix.\\
\vspace{.4cm}
Fourth, multiply all of them to produce an $(n \times 1)$ vector.\\
\vspace{.4cm}
Use forumula to sum the product of each component with each data point to produce a new $(d \times 1)$ vector.\\
\vspace{.4cm}
The optimal f(x) is given by $f^*(x)=\beta^{\top}x$ with $\beta^{\top}$ given above.













































\end{flushleft}
\end{document}